% Source-transform material can follow the Euclidean paragraph at the end of
% canonical_local.tex (existing source-kernel/amplitude-interface section).
% The final two subsections can instead be placed after modes_predecessor.tex
% inside Part D. All headings below are subsections, preserving top-level count.
\subsection{An exact transform of the original continuum}
\label{subsec:actual-continuum-transform-update}
The owner's review derives an elementary transform of the same normalized
source measure. Write
\[
 F_E(\tau)=F(\tau/R),\qquad
 F(a)=\frac{10}{\pi}\int_0^\infty e^{-ax}j_2(x)^2\,dx.
\]
The Legendre representation gives the exact overlap
\begin{align*}
 C(v)&=\int_{-1}^{1-v}P_2(t)P_2(t+v)\,dt\\
 &=\frac25-v+\frac{v^3}{2}-\frac{3v^5}{40},
       \qquad 0\leq v\leq2,\\
 F(a)&=\frac5\pi\int_0^2 C(v)\frac{a}{a^2+v^2}\,dv.
\end{align*}
Direct integration therefore yields
\begin{align}
 16\pi F(a)&=12a^3+56a+32\arctan(2/a)\nonumber\\
 &\quad-a(3a^4+20a^2+40)\log(1+4/a^2).
 \label{eq:exact-continuum-transform-update}
\end{align}
The original positive Bessel integral fixes the sign; $C$ itself is signed.
Its first cancellations, $\int C=\int v^2C=0$, explain why ordinary
floating-point subtraction loses digits at large $a$.

Range-reduced rational logarithm and arctangent series, with explicit
remainders and a Machin enclosure of $\pi$, give outward intervals of
width below $10^{-75}$ at all operands in the following table. Entries
are midpoint displays; exact rational endpoints are in the source packet.
\[
 \begin{array}{c|c}
 a&F(a)\\\hline
 0.01&0.9236282572354819\\
 0.1&0.6004698962807585\\
 1&0.040468535993562725\\
 10&3.254795554745759\times10^{-6}\\
 100&3.3938508389223904\times10^{-11}
 \end{array}
\]
A convergent large-$a$ alternative uses
\[
 \int_0^2v^mC(v)\,dv
 =\frac{2^{m+2}m(m-2)}{(m+1)(m+2)(m+4)(m+6)}.
\]
After $N$ geometric terms its absolute transform remainder is at most
$4(4/a^2)^N/[\pi a(2N+1)]$, for $a>2$. The small-$a$ expansion is
\[
 \begin{aligned}
 F(a)=1&+\frac5\pi a\log a\\
       &+\frac5\pi\left(\frac12-\log2\right)a
         +O(a^3|\log a|).
 \end{aligned}
\]
Its infinite initial right slope and late-time power-law tail both retain
the continuous, gapless source rather than a finite family of rates.

\subsection{Soft-spectrum bounds without a minimum rate}
\label{subsec:source-soft-spectrum-update}
The compact Bessel representation implies $|j_2(x)|\leq x^2/15$.
For the original $d\nu_R(k)=(10R/\pi)j_2(kR)^2\,dk$, this gives
\begin{align}
 F_E(\tau)&\leq\min\!\left(1,\frac{16R^5}{15\pi\tau^5}\right),\nonumber\\
 \nu_R([0,q])&\leq\min\!\left(1,\frac{2(Rq)^5}{225\pi}\right).
 \label{eq:gapfree-source-bounds-update}
\end{align}
Suppose a physical reduction independently establishes
$\Xi(s)=A+(1-A)F(\lambda s/R)$ with $0\leq A\leq1$ and
$\lambda\geq\lambda_{\min}>0$. Set $E(s)=F(\lambda_{\min}s/R)$. Then
\[
 \max\!\left(0,\frac{\Xi(s)-E(s)}{1-E(s)}\right)
 \leq A\leq\Xi(s),\qquad s>0.
\]
Simultaneous data bands and outward upper $E$ values can be mapped and
intersected. The conversion-scale restriction is distinct from a hard
spectral gap. The observable and its real-time or logarithmic-redshift
identification still require the physical reduction.

More generally let $\nu([0,p])\leq Cp^5$, let the transient amplitude
measure be dominated by $B\nu$, and let an omitted submeasure obey
$\int(1-e^{-ph})\,d\beta_{\rm omitted}\leq m$, with $h>0$.
Splitting at any $p_0>0$ gives
\begin{align*}
 B_{\rm omitted}&\leq B,\\
 B_{\rm omitted}&\leq BCp_0^5+\frac{m}{1-e^{-p_0h}}.
\end{align*}
Using $(1-e^{-x})^{-1}\leq1+x^{-1}$ and optimizing at
$p_0=[m/(5BCh)]^{1/6}$ yields
\begin{equation}
 \begin{gathered}
 B_{\rm omitted}\leq B,\\
 B_{\rm omitted}\leq m+6(BC)^{1/6}
                          \left(\frac{m}{5h}\right)^{5/6}.
 \end{gathered}
 \label{eq:soft-spectrum-payment-update}
\end{equation}
Zero-mass cases follow by a limit. For the mapped original source,
$C=2(R/\lambda_{\min})^5/(225\pi)$. Thus a polynomial residual bound
$m\leq\eta/d^2$ can supply a plateau payment after these source-class
operands are specified, without imposing a nonexistent minimum continuum
rate.

% Inference insertion boundary: this material can be moved intact to Part D.
\subsection{Infinite exact rank and positive Gaussian surrogates}
\label{subsec:continuum-Gauss-update}
On a Euclidean grid $\tau_j=\tau_0+jh$, map $k$ to
$\lambda=e^{-hk}$. The actual difference measure is
\[
 d\mu=B e^{-\tau_0k}(1-e^{-hk})d\nu_R(k).
\]
For $B>0$, it has continuous support in $(0,1)$ apart from isolated Bessel
zeros. A nonzero polynomial has finitely many roots, so $\int|p|^2d\mu>0$.
Every finite Hankel matrix is positive definite. The original continuum
therefore has infinite exact Hankel rank.

Let $p_r$ be its monic orthogonal polynomial. Its $r$-node Gaussian
quadrature has positive weights and is exact through degree $2r-1$.
Division by $p_r$ establishes exactness; integrating the squared cardinal
polynomials establishes positivity of the weights. Anchor the surrogate
at the actual $Y_0$, and put $q_j(\lambda)=\sum_{k=0}^{j-1}\lambda^k$.
Hermite interpolation then gives
\begin{align}
 \widetilde Y_j&=Y_j,\qquad 0\leq j\leq2r,\nonumber\\
 0\leq\widetilde Y_j-Y_j
 &\leq\binom{j}{2r+1}\int p_r^2\,d\mu,
       \qquad j\geq2r+1.
 \label{eq:continuum-Gauss-holdout-update}
\end{align}
Here $q_j^{(2r)}/(2r)!$ is nonnegative and bounded by the binomial
coefficient on $[0,1]$; the first missed error equals $\int p_r^2d\mu$.
If the true amplitude integral is finite and its plateau is nonnegative,
monotone convergence of $q_j$ also makes the surrogate plateau no smaller
than the true plateau.

\subsection{Certified surrogates of the actual source curve}
For $Y_j=F(j)$, namely $\tau_0=0$ and $h=R$, the true plateau is zero.
Full interval arithmetic gives the following nodes and difference weights;
these are midpoint displays of certified source intervals.
\[
 \begin{array}{c|cc}
 r&\lambda_i&w_i\\\hline
 1&0.03729697269&0.95953146401\\
 2&0.02123986980&0.87066215903\\
  &0.19461014031&0.08886930497
 \end{array}
\]
The corresponding plateaus and first missed errors are
\[
 \begin{array}{c|cc}
 r&\widetilde A&\widetilde Y_{2r+1}-Y_{2r+1}\\\hline
 1&0.003294435783&0.002423770451\\
 2&0.0001005695591&0.00003902006495
 \end{array}
\]
The exact $H_2$ and $H_3$ determinant enclosures are strictly positive,
near $0.002325684008961166$ and $9.074834108661564\times10^{-8}$.
The packet retains all moment, node, weight, plateau and signed holdout
intervals through $j=7$. These are controlled finite approximations to the
actual continuous Euclidean source. Their surplus plateaus illustrate why
an exact fit on a short grid is distinct from a physical finite spectrum.
Noisy mode recovery retains the chapter's simultaneous feasibility and
covariance treatment.
